%Hoja de vida Rodrigo Taborda
%Curriculum Vitae Rodrigo Taborda

\documentclass[12pt,a4paper,helvetica,narrow,nologo,notitle,booktabs,spanish]{europecv}
\usepackage{graphicx}

\setlength{\voffset}{-1in}          % Reset the top margin back to zero
\setlength{\hoffset}{-1in}          % Reset the left margin to zero
\setlength{\topmargin}{2cm}       % Top margin
\setlength{\oddsidemargin}{2cm}   % Odd side left margin to 30mm
\setlength{\evensidemargin}{2cm}  % Even side left margin to 30mm
\setlength{\textwidth}{16.9cm}  % Text width to be 16.9cm (an A4 page
                                % is 16.9cm+2cm+2cm=21.5cm wide)
\setlength{\textheight}{24.7cm} % Text height to be 24.7cm (an A4
                                % page is 21.7cm+2cm+2cm=29.7cm long)
%\usepackage[latin1]{inputenc}

%Section of personal information
    \ecvname{Taborda, Rodrigo}
    \ecvfootername{Rodrigo Taborda}
    \ecvaddress{
    Calle 21 \# 1 - 20. Facultad de Administración, Universidad de Los Andes. Bogotá, Colombia
    }
    \ecvemail{\emph{%    rodrigo.taborda@anu.edu.au,
                    %    rotabori@gmail.com,
                         r.taborda@uniandes.edu.co,
                         www.rodrigotaborda.com}}
    \ecvnationality{Colombia}
%    \ecvdateofbirth{Agosto 26, 1975}
%    \ecvpicture[height=5cm]{cv_pic_6.eps}
%        \ecvbeforepicture{\raggedleft}

\begin{document}
%Some refinements
    \fancyfoot{}
    \fancyfoot[R]{\thepage}

\begin{europecv}

    \ecvpersonalinfo %Introduces personal information included above

%---------------------------------------------------------------------
    \cline{2-2}
    \ecvsection{Cargos}
        \ecvitem*{}{2018 - Actualidad. \textbf{Profesor asociado}. Facultad de Administración, Universidad de los Andes.}

        \ecvitem*{}{2013 - 2018. \textbf{Profesor asistente}. Facultad de Administración, Universidad de los Andes.}

        \ecvitem*{}{2010 - 2013. \textbf{Profesor}. Facultad de Economía, Universidad del Rosario.}

        \ecvitem*{}{2010 - 2023. \textbf{Editor}. \emph{Revista de Economía del Rosario}.}

        \ecvitem*{}{2008 - 2010. \textbf{Consultor}. ACIL-Tasman, Canberra, Australia.}

        \ecvitem*{}{2002 - 2004. \textbf{Profesor}. Universidad del Rosario.}

%---------------------------------------------------------------------
    \cline{2-2}
    \ecvsection{Educación}

        \ecvitem{Pregrado}{}

            \ecvitem*{}{\textbf{Economía}. Universidad de Antioquia, Medellín-Colombia. 1993-2000.
        Artículo de grado: ``Inconsistencia temporal de políticas óptimas, reglas y política anti-inflacionaria en Colombia (1963-1998)''}

            \ecvitem*{}{\textbf{Cursos de pregrado en economía}. Indiana University of Pennsylvania. Indiana, PA. USA. 1998.}

        \ecvitem{Postgrado}{}

            \ecvitem*{}{\textbf{Ph.D. Economics}. Research School of Asian Pacific-RSPAS- Economics. Australian National University, Canberra-Australia. Tesis: ``Manufacturing industry performance after structural reforms. The case of Colombia.'' Supervisor: Prema-Chandra Athukorala. Diciembre 2009.}

            \ecvitem*{}{\textbf{M.A. Economics}. Australian National University, Canberra-Australia. 2005}

            \ecvitem*{}{\textbf{Maestría Economía}. Universidad del Rosario, Bogota-Colombia. 2001-2002. (Título otorgado en Agosto de 2004). Tésis de grado: ``Technical efficiency and performance in the electric energy distribution system of Colombia''}

        \ecvitem{Cursos y seminarios}{}

            \ecvitem*{}{\textbf{Topics in empirical finance}. Marno Verbeek (Erasmus University). Mayo 27-31, 2013. Departamento de economía. Universidad del Rosario.}

            \ecvitem*{}{\textbf{Analysis of microdata}. Rainer Winkelmann (University of Zurich). Febrero 11-15, 2013. Departamento de economía. Universidad del Rosario.}

            \ecvitem*{}{\textbf{Time series for macroeconomics and finance}. Andrew Harvey (Cambridge University). Marzo 12-16, 2012. Departamento de economía. Universidad del Rosario.}

            \ecvitem*{}{\textbf{Econometric analysis of panel data}. Badi Baltagi (Syracuse University). Oct 24-28, 2011. Departamento de economía. Universidad del Rosario.}

            \ecvitem*{}{\textbf{Econometric analysis of panel data}. Badi Baltagi (Texas A\&M University). Mayo 26-30, 2003. Departamento de economía. Universidad del Rosario, Bogotá.}

            \ecvitem*{}{\textbf{Non-linear models in economics and finance}. Phillip H. Franses (Erasmus University). Noviembre 14-18, 2000. Departamento de economía. Universidad del Rosario, Bogotá.}

%---------------------------------------------------------------------
    \cline{2-2}
    \ecvsection{Publicaciones}

        \ecvitem{Revistas\\
        \footnotesize{\emph{Artículos en revistas académicas}}}{}

        \ecvitem*{}{``\textbf{A Test of Auditor Performance under the Service Paradigm: Do Emotional Intelligence and National Culture Add to Our Understanding?}'' (Con Rosa González, Stephen B. Salter y Felipe Janica). Journal of International Accounting, Auditing and Taxation. En prensa (2027).}

        \ecvitem*{}{``\textbf{Time Poverty: An Unintended Consequence of Women Participation in Farmers' Associations}'' (With Andres Barrios, Ximena Rueda). Social Indicators Research: An International and Interdisciplinary Journal for Quality-of-Life Measurement. Springer, vol. 178(1), pages 225-253, 2025.}

        \ecvitem*{}{``\textbf{Fund Flows, Stock Markets, and Economic Policy Uncertainty: From the Perspective of a CIVETS Nation}'' (With Joseph J French, Philipp D Schaberl). Emerging Markets Finance and Trade. Vol. 60 No. 5. pages 967-982. 2023.}

        \ecvitem*{}{``\textbf{Digitalization and e-government in the lives of urban migrants: evidence from Bogota}'' (With Charles Martin, Sonia Camacho, Constantine Ruhe). Policy \& Internet. Vol. 14 No. 2. 2022.}

        \ecvitem*{}{``\textbf{Savings and Gender in Colombia: ANOVA, Simple Linear Regression, and Extreme Observations}'' (With Luz Orozco). Sage research cases 2022.}

        \ecvitem*{}{``\textbf{How does the training of subsistence entrepreneurs affect their hope levels and perceived goal-attainment?}'' (With Andr\'es Barrios y Ezequiel Reficco). Education and training. Vol. 61 No. 7/8. 2019.}

        \ecvitem*{}{``\textbf{Family firms and financial performance. A meta-regression analysis}'' (With Maximiliano Gonz\'alez y Juan Idrobo). Academia - Latin American Journal. Vol. 32. No. 3. 2019}

        \ecvitem*{}{``\textbf{Disentangling the relationship between liquidity and returns in Latin America}'' (With Joseph French). Global finance journal. Vol. 36. 2018}

        \ecvitem*{}{``\textbf{Stock liquidity and second blockholder as drivers of corporate
            value: Evidence from Latin America}'' (With Carlos Pombo). International review of economics and finance. Vol. 51. 2017}

        \ecvitem*{}{``\textbf{Innovation and Productivity in the Colombian Service and Manufacturing industries}'' (With Juan Gallego y Luis Gutiérrez). Emerging Markets Finance and Trade. Vol. 51. 2015.}

        \ecvitem*{}{``\textbf{Procedural transparency in Latin American central banks under inflation targeting schemes. A text analysis of the minutes of the Boards of Directors}''. Ensayos sobre Política económica, Banco de la República. Vol. 33. 2015.}

        \ecvitem*{}{``\textbf{Insurance use and economic growth in Latin America. Panel data evidence}'' (With Ángela Concha). Lecturas de Economía. Vol. 81. Julio - Diciembre 2014.}

        \ecvitem*{}{``\textbf{Bias in economic news: the reporting of the nominal exchange rate behavior in Colombia}''. Journal Economia - LACEA. Vol. 14(1). Fall. 2013.}

        \ecvitem*{}{``\textbf{Manufacturing employment and wage differentials after structural adjustment reforms in Colombia}'' (Con Juan Guataquí). Lecturas de Economía. Vol. 48. Enero - Junio 2011.}

        \ecvitem*{}{``\textbf{Ownership and control in Colombian corporations}'' (Con Carlos Pombo and Luis Gutierrez). Quarterly Review of Economics and Finance. Vol. 48(1) 2008.}

        \ecvitem*{}{``\textbf{Technical instruments and conceptual definitions: theoretical and empirical implications of the change on the definition of unemployment in Colombia}'' (Con Juan Guataqui). Revista de Economía del Rosario. vol. 9(1) June 2006.}

        \ecvitem*{}{``\textbf{Performance, efficiency and productivity in Colombia's power distribution system: effects of the 1994 regulatory reform}''. (Con Carlos Pombo). Energy Economics. Vol. 28(3) 2006.}

        \ecvitem{Libros}

        \ecvitem*{}{``Ejercicios de econometría''. Editorial Universidad del Rosario. 2016}

        \ecvitem*{}{``Curso de macroeconomía''. Editorial Universidad del Rosario. 2005}

        \ecvitem{Revisión de libros}

        \ecvitem*{}{``Playing Nostradamus with Latin America. A review of \emph{Latin America's political economy of the possible}'' por Javier Santiso (MIT 2006). Agenda, Vol. 14, No. 1. 2007.}

        \ecvitem{Documentos de trabajo \\
        \footnotesize{\emph{pre-publicaciones académicas}}}{}

        \ecvitem*{}{``Stock liquidity and second blockholder as drivers of corporate value: Evidence from Latin America''. Documento del CEDE No. 014074. Facultad de Economía. Universidad de los Andes. Bogotá.}

        \ecvitem*{}{``Innovation and Productivity in the Colombian Service Industry''. Inter-American Development Bank working paper series. No. 81079. IADB. Washington.}

        \ecvitem*{}{``Media bias and central bank response. Evidence from the nominal exchange rate behavior in Colombia''. Documento de trabajo No. 115. Facultad de Economía. Universidad del Rosario. Bogota.}

        \ecvitem*{}{``Productivity growth in electric energy retail in Colombia. A bootstrapped Malmquist indices approach''. Documento de trabajo No. 114. Facultad de Economía. Universidad del Rosario. Bogota.}

        \ecvitem*{}{``Ownership and control in Colombian corporations''. (Documento de trabajo No. 66.)
            Borradores de Investigación No. 66. Facultad de Economía. Universidad del Rosario. Bogota.
            También en IADB-RES Documento de trabajo.}

        \ecvitem*{}{``Wage rigidity and turnover costs in Colombian manufacturing industry''. (Documento de trabajo No. 37.) Borradores de Investigación No. 37. Facultad de Economía. Universidad del Rosario. Bogota.}

        \ecvitem*{}{``Technical instruments and conceptual definitions: theoretical and practical
            implications of the change on the definition of unemployment in Colombia''. (Documento de trabajo No. 36.) Borradores de Investigación No. 36. Facultad de Economía. Universidad del Rosario. Bogota.}

        \ecvitem*{}{``Estudiar economía y racionalidad: un experimento''. (Documento de trabajo No. 34.) Borradores de Investigación No. 34. Facultad de Economía. Universidad del Rosario. Bogota.}

        \ecvitem*{}{``Política monetaria y cambiaria.'' En: Perfil de Coyuntura Económica. Diciembre 1999. Centro de Investigaciones Económicas (CIE). Universidad de Antioquia. Medellín.}

        \ecvitem*{}{``Sector Agropecuario.'' En: Perfil de Coyuntura Económica. Junio 1998. Centro de Investigaciones Económicas (CIE). Universidad de Antioquia. Medellín.}

        \ecvitem*{}{``Indicadores fiscales.'' En: Lecturas de Economía. Julio - Diciembre 1997. Centro de Investigaciones Económicas (CIE). Universidad de Antioquia. Medellín.}

        \ecvitem*{}{``Indicadores de salud.'' En: Lecturas de Economía. January - Julio 1997. Centro de Investigaciones Económicas (CIE). Universidad de Antioquia. Medellín.}

        \ecvitem{Traducciones (inglés--español)}{}
            \ecvitem*{}{2026. \textbf{Oportunismo y buena voluntad. Expansión empresarial canadiense en Colombia, 1867--1979}. Editorial Universidad del Rosario, en prensa. Original: \emph{Opportunism and goodwill: Canadian business expansion in Colombia, 1867--1979}, Stefano Tijerina, University of Toronto Press, 2021.}

            \ecvitem*{}{2026. \textbf{¿Quién gobierna a los excluidos?} Editorial Universidad del Rosario. Original: \emph{Governing the excluded: rural livelihoods beyond coca in Colombia's peace laboratory}, Alex Diamond, University of Chicago Press, 2026.}

            \ecvitem*{}{2024. \textbf{Género y ciudadanía en justicia transicional. La cotidianidad de la reparación y reintegración en Colombia}. Universidad del Rosario. Original: \emph{Gender and citizenship in transitional justice. Everyday experiences of reparation and reintegration in Colombia}, Sanne Weber, Bristol University Press.}

            \ecvitem*{}{2021. \textbf{Las reformas agrarias en Latinoamérica: restricciones institucionales y división de las élites}. Universidad del Rosario. Original: \emph{Autocracy and redistribution. The politics of land reform}, Michael Albertus, Cambridge University Press.}

            \ecvitem*{}{2015. \textbf{The development of the agricultural sector of Colombia into the 1960s}. Albert Berry (Editorial Universidad del Rosario, 2017, edición en inglés). Traducción al español para la Universidad del Rosario.}

            \ecvitem*{}{2013. \textbf{Demasiado lejos para exportar}, capítulo sobre Brasil (con María José Molano). Banco Interamericano de Desarrollo. Original: \emph{Too far to export: domestic transport costs and regional export disparities in Latin America and the Caribbean}, BID.}

            \ecvitem*{}{2008. \textbf{Sistemas fiscales intergubernamentales}. Mayol S.A. Original: \emph{Intergovernmental fiscal transfers. Principles and practice}, Robin Boadway y Anwar Shah (eds.), Banco Mundial.}

            \ecvitem*{}{2007. \textbf{Sistemas y estándares de calidad hacia la construcción de ventaja competitiva}. Mayol S.A. Original: \emph{Quality systems and standards for a competitive edge}, Banco Mundial.}

            \ecvitem*{}{2006. \textbf{Estudio del financiamiento de la salud: una guía para especialistas}. Mayol S.A. Original: \emph{Health financing revisited. A practitioner's guide}, Banco Mundial.}

            \ecvitem*{}{2006. \textbf{Economic growth in the 1990s: learning from a decade of reform}. Banco Mundial. Traducción al español.}

            \ecvitem*{}{2005. \textbf{Lecciones del TLCAN para América Latina y el Caribe}. Alfaomega S.A. Original: \emph{Lessons from NAFTA for Latin America and the Caribbean}, Stanford University Press -- Banco Mundial.}

            \ecvitem*{}{2004. \textbf{Pobreza en Colombia: tendencias divergentes en desarrollo económico y social}. Banco Mundial. Original: \emph{Poverty in Colombia: the divergent trends of economic and social development}.}

            \ecvitem*{}{2004. \textbf{Inequality in Latin America: breaking with history?} Banco Mundial. Traducción al español.}

            \ecvitem*{}{2004. \textbf{Bienes públicos internacionales: incentivos, medición y financiación}. Alfaomega S.A. y Banco Mundial. Original: \emph{International public goods: incentives, measurement and financing}, Kluwer Academic Publishers -- Banco Mundial.}

            \ecvitem*{}{2003. \textbf{Regional integration and development}. Maurice Schiff y L. Alan Winters, Banco Mundial. Traducción al español para Alfaomega S.A.}

            \ecvitem*{}{2003. \textbf{The limits of stabilization: infrastructure, public deficits, and growth in Latin America}. William Easterly y Luis Servén (eds.), Banco Mundial. Traducción al español.}

            \ecvitem*{}{2002. \textbf{Globalización, crecimiento y pobreza}. Alfaomega S.A. y Banco Mundial. Original: \emph{Globalization, growth and poverty}, Banco Mundial.}

            \ecvitem*{}{2002. \textbf{Fronteras de la economía del desarrollo}. Alfaomega S.A. Original: \emph{Frontiers of development economics. The future in perspective}, Joseph Stiglitz y Gerald Meier (eds.), Oxford University Press.}

        \ecvitem{Revisión técnica de traducciones}{}
            \ecvitem*{}{2012. \textbf{Construyendo el futuro de la relación entre Asia y el Pacífico y América Latina y el Caribe}. Banco Interamericano de Desarrollo. Revisión de la traducción.}

            \ecvitem*{}{2003 - 2004. \textbf{Informe sobre el desarrollo mundial 2003 y 2004}. Banco Mundial y Alfaomega S.A. Revisión técnica de la traducción.}

            \ecvitem*{}{2003. \textbf{Cerrar la brecha en educación y tecnología}. Banco Mundial y Alfaomega S.A. Revisión técnica de la traducción.}

%---------------------------------------------------------------------
    \cline{2-2}
    \ecvsection{Conferencias}

        \ecvitem*{}{``Family firms and financial performance. A meta-regression analysis''. Universidad de los Andes, Oct. 2017. EAFIT, Jun. 2018. Universidad Javeriana, Ago. 2018.}

        \ecvitem*{}{``Stock liquidity and second blockholders as drivers of corporate value: evidence from Latin America''. World Finance Conference. Julio 26-28. 2016.}

        \ecvitem*{}{``Procedural transparency in Latin America’s central banks. A text analysis of the boards’ minutes''. XI Seminario de la Edición Especial de la Revista ESPE Cambios en la Política Monetaria y en la Banca Central en las últimas dos décadas, Bogotá, Colombia. October 23-24. 2014.}

        \ecvitem*{}{``Automatic generation of a personalized answer to a problem set''. Stata Conference, New Orleans, USA. July 18-19. 2013.}

        \ecvitem*{}{``Combinación del uso de Stata y \LaTeX en la evaluación de un curso de econometría o estadística''. Encuentro Internacional de usuarios de Stata, Mexico D.F., Mexico. May 3. 2013.}

        \ecvitem*{}{``Productivity growth in electric energy retail in Colombia. A bootstrap approach''. 51st Meeting of the EWGCFM / 1st Conference of the RCEM \& ICSTF, London, UK. May 16-18. 2013.}

        \ecvitem*{}{``Media bias and central bank response. Evidence from the nominal exchange rate behavior in Colombia''. LACEA Panel - Economía Journal. Washington, DC, USA, Abril 13-14. 2012.}

        \ecvitem*{}{``Productivity growth in electric energy retail in Colombia. A bootstrap approach''. 30th United States Association for Energy Economics / International Association for Energy Economics (USAEE / IAEE) North American Conference, Washington, DC, USA. Octubre 9-12. 2011.}

        \ecvitem*{}{``Media bias and central bank response. Evidence from the nominal exchange rate behavior in Colombia''. Lupcon Center for Business Research. Frankfurt, Germany. Julio 5-6. 2011.}

        \ecvitem*{}{``Measuring media bias on the economy: an application to the behavior of the nominal exchange rate in Colombia''. Latin American and Caribbean Economic Association (LACEA). Medellín. Noviembre 11-13. 2010.}

        \ecvitem*{}{``Effect of import tariff reduction on manufacturing productivity in Colombia''.  Queensland University of Technology - Australian National University PhD workshop 2008 PhD Conference in Economics and Business. Australian National University. Canberra. Noviembre 20-21. 2008.}

        \ecvitem*{}{``Structural adjustment reforms and productivity heterogeneity: Evidence from the Colombian manufacturing sector''.  2008 PhD Conference in Economics and Business. Australian National University. Canberra. Noviembre 20-21. 2008.}

        \ecvitem*{}{``Performance and efficiency of Colombia's Energy distribution system''. 23 International Association for Energy Economics. North American Conference. Mexico. Mexico. Octubre 19-21. 2003.}

        \ecvitem*{}{``Regulatory reform and electric energy demand in the Colombian Industry'' (``Reforma        regulatoria y demanda de energía eléctrica en la industria colombiana''). 5th Latin American Congress: Electricity Generation and Transmission. Sao Pedro. Brasil. Noviembre 2003.}

%---------------------------------------------------------------------
    \cline{2-2}

% (sección comentada: no relevante por ahora)
%     \ecvsection{Investigación}
%     \ecvitem{\footnotesize{\emph{Participación en los siguientes proyectos de investigación}}}{}

%         \ecvitem*{}{\textbf{Manufacturing industry electricity demand. A time series approach.} Enero 2013 - Diciembre 2013. Grupo de trabajo: Eduardo Soto. Universidad del Rosario, Colombia.}
 	
%         \ecvitem*{}{\textbf{Insurance and economic growth in Latin America. Panel data evidence.} Enero 2012 - Diciembre 2012. Grupo de trabajo: Angela Concha. Universidad del Rosario, Colombia.}

%         \ecvitem*{}{\textbf{Text analysis of three Latin American central banks' minutes.} Junio 2012 - Junio 2013. Universidad del Rosario, Colombia.}

%         \ecvitem*{}{\textbf{Internet and wages. Evidence from Colombia.} Enero 2012 - Junio 2013. Grupo de trabajo: Juan Oviedo, Eduardo Soto y Juan Guataqui. Universidad del Rosario, Colombia.}

%         \ecvitem*{}{\textbf{Productivity in electric energy retail service in Colombia. A bootstrapped malmquist indices approach.} Jan 2012 - June 2013. Universidad del Rosario, Colombia.}

%         \ecvitem*{}{\textbf{Estimación de productividad en servicio de comercialización de energía eléctrica en Colombia.} Octubre 2010 - Junio 2011. Grupo de trabajo: Manuel Ramírez y Luis H. Gutiérrez. Universidad del Rosario, Colombia.}

%         \ecvitem*{}{\textbf{Voting behavior India.} Julio 2009 - Agosto 2009. Asistente de investigación bajo la supervision de Raghbendra Jha. Australian National University, Australia.}

%         \ecvitem*{}{\textbf{Economic Growth, Law and Corruption: Evidence from India.} Mayo 2009 - Julio 2009. Asistente de investigación bajo la supervision de Raghbendra Jha. Australian National University, Australia.}

%         \ecvitem*{}{\textbf{Economic growth in Malaysia.} Octubre 2008 - Diciembre 2008. Asistente de investigación bajo la supervision de Prema-Chandra Athukorala. Australian National University, Australia.}

%         \ecvitem*{}{\textbf{Poverty in Thailand.} Junio 2008 - Diciembre 2008. Asistente de investigación bajo la supervision de Peter Warr. Australian National University, Australia.}

%         \ecvitem*{}{\textbf{Burma economy's structure.} Diciembre 2007 - Junio 2008. Asistente de investigación bajo la supervision de Peter Warr. Australian National University, Australia.}

%         \ecvitem*{}{\textbf{Fair trade.} Junio - Julio 2007. Asistente de investigación bajo la supervision de Martin Richardson. Australian National University, Australia.}

%         \ecvitem*{}{\textbf{Local content schemes in broadcasting. The Australian experience.} Diciembre 2006 - Mayo 2007. Asistente de investigación bajo la supervision de Martin Richardson. Australian National University, Australia.}

%         \ecvitem*{}{\textbf{Corporate governance in Latin America and the Caribbean}. Inter-American development Bank (IADB). 2003-2004. Grupo de trabajo: Carlos Pombo, Luis Gutierrez and Rodrigo Taborda. Reporte final: ``Ownership and Control in Colombian Corporations'' con Carlos Pombo and Luis Gutierrez; publicado en Quarterly Review of Economics and Finance. Universidad del Rosario,        Colombia.}

%         \ecvitem*{}{\textbf{Performance and efficiency in thermal generation and urban electricity distribution of energy in Colombia: effects of the 1994 regulatory reform}. COLCIENCIAS 2003. Grupo de trabajo: Rodrigo Taborda and Carlos Pombo. Documento final publicado en Energy Economics. Universidad del Rosario, Colombia.}

%         \ecvitem*{}{\textbf{Costs and benefits of privatization in Latin America}. Inter-American Development Bank (IADB). 2002-2003. Documento final: ``Privatization in Colombia: a plant performance analysis''. Documento final publicado como capitulo en ``Privatization in Latin America Myths and Reality'' Stanford University Press. Asistente de investigación bajo la supervisión de Carlos Pombo y Manuel Ramirez. Universidad del Rosario, Colombia.}

%         \ecvitem*{}{\textbf{Regulatory reform in Colombia's electric utilities}. Asistente de investigación bajo la supervisión de Carlos Pombo. Documento final publicado en Quarterly Review of Economics and Finance. Universidad del Rosario, Colombia.}

    \ecvsection{Docencia}
        \ecvitem*{}{\begin{itemize}\item 29 cursos distintos dictados entre 2002 y 2026 en cuatro instituciones: 19 en la Universidad de los Andes y 8 en la Universidad del Rosario (Colombia), y 2 fuera de Colombia (IESA, Venezuela; Universidad Americana, Paraguay), además de tutorías en la Australian National University. \item 14 de pregrado, 14 de maestría y MBA y 1 de doctorado, más 2 programas de educación ejecutiva. \item 10 cursos de diseño original y 4 en línea (Maestría en Analítica y Gestión Financiera y MBA virtual). \item Al menos 133 cursos-semestre en total, de los cuales 16 intersemestrales intensivos. \item Docencia en español e inglés (Maestría Internacional en Finanzas con Nova SBE) y en codocencia. \end{itemize}}

        \ecvitem{Universidad de los Andes -- Pregrado}{}
            \ecvitem*{}{2013 - Actualidad. \textbf{Análisis del entorno colombiano}. Pregrado en Administración. Periodos (18 semestres): 2013-2, 2014-1, 2014-2, 2015-1, 2015-int, 2015-2, 2018-1, 2018-2, 2019-1, 2019-2, 2020-1, 2021-2, 2022-1, 2023-1, 2024-1, 2024-2, 2025-1, 2026-2.}

            \ecvitem*{}{2017 - Actualidad. \textbf{Herramientas para la toma de decisiones I (análisis multivariado; Modelos multivariados)}. Pregrado en Administración. Periodos (22 semestres): 2017-1, 2017-2, 2018-1, 2018-2, 2019-1, 2019-int, 2019-2, 2020-1, 2020-int, 2020-2, 2021-1, 2021-2, 2022-1, 2022-2, 2023-1, 2023-2, 2024-1, 2024-2, 2025-1, 2025-int, 2025-2, 2026-int. Diseño original del curso.}

            \ecvitem*{}{2025 - Actualidad. \textbf{Visualización de datos}. Pregrado en Administración. Periodos (4 semestres): 2025-1, 2025-2, 2026-1, 2026-2. Diseño original del curso.}

            \ecvitem*{}{2026 - Actualidad. \textbf{Introducción a la ciencia de datos}. Pregrado en Ciencia de Datos. Periodos (2 semestres): 2026-1, 2026-2.}

            \ecvitem*{}{2016 - 2020. \textbf{Economía gerencial}. Pregrado en Contaduría Internacional. Periodos (9 semestres): 2016-1, 2016-2, 2017-1, 2017-2, 2018-1, 2019-1, 2019-2, 2020-1, 2020-2. Diseño original del curso.}

            \ecvitem*{}{2018. \textbf{Planeación financiera (con Maximiliano González)}. Pregrado en Administración. Periodos (1 semestre): 2018-2.}

        \ecvitem{Universidad de los Andes -- Posgrado}{}
            \ecvitem*{}{2015 - Actualidad. \textbf{Análisis multivariado (Métodos I / Diseño de investigación I)}. Doctorado en Administración y Maestría en Investigación en Administración. Periodos (11 semestres): 2015-1, 2015-2, 2016-1, 2017-1, 2017-2, 2018-1, 2019-1, 2020-1, 2021-1, 2025-1, 2026-1. Diseño original del curso.}

            \ecvitem*{}{2017 - Actualidad. \textbf{Microeconomía financiera}. Maestría en Finanzas. Periodos (18 semestres): 2017-2, 2018-2, 2019-1, 2019-2, 2020-1, 2020-2, 2021-1, 2021-2, 2022-1, 2022-2, 2023-1, 2023-2, 2024-1, 2024-2, 2025-1, 2025-2, 2026-1, 2026-2.}

            \ecvitem*{}{2014 - 2022. \textbf{Econometría para finanzas (Métodos de regresión para finanzas)}. Maestría en Finanzas. Periodos (6 semestres): 2014-1, 2018-2, 2019-int, 2020-int, 2021-int, 2022-int.}

            \ecvitem*{}{2018 - 2023. \textbf{Intermediate econometrics for finance (Financial econometrics)}. Maestría Internacional en Finanzas, Universidad de los Andes--Nova School of Business and Economics. Periodos (6 semestres): 2018-2, 2019-2, 2020-2, 2021-2, 2022-2, 2023-2. Diseño original del curso.}

            \ecvitem*{}{2014. \textbf{Evaluación de impacto}. Maestría en Gerencia y Práctica del Desarrollo (MDP). Periodos (1 semestre): 2014-1.}

            \ecvitem*{}{2025. \textbf{Econometría intermedia (Métodos cuantitativos)}. Maestría en Analítica y Gestión Financiera (virtual). Periodos (2 semestres): 2025-1, 2025-2. Diseño original del curso.}

            \ecvitem*{}{2026 - Actualidad. \textbf{Econometría financiera (Econometría avanzada para finanzas)}. Maestría en Analítica y Gestión Financiera (virtual). Periodos (1 semestre): 2026-int. Diseño original del curso.}

            \ecvitem*{}{2026 - Actualidad. \textbf{Simulación de escenarios y riesgo}. Maestría en Analítica y Gestión Financiera (virtual). Periodos (1 semestre): 2026-2.}

            \ecvitem*{}{2023 - Actualidad. \textbf{Estadística (Incertidumbre y análisis de datos)}. MBA. Periodos (4 semestres): 2023-2, 2024-2, 2025-2, 2026-2.}

            \ecvitem*{}{2024 - Actualidad. \textbf{Análisis y visualización de datos}. MBA Ejecutivo. Periodos (2 semestres): 2024-int, 2026-int. Diseño original del curso.}

            \ecvitem*{}{2023 - Actualidad. \textbf{Blockchain para la empresa}. MBA (virtual). Periodos (4 semestres): 2023-int, 2024-int, 2025-int, 2026-int. Diseño original del curso.}

            \ecvitem*{}{2024 - 2025. \textbf{Análisis de datos para mercadeo (Marketing analytics)}. Maestría en Mercadeo. Periodos (2 semestres): 2024-1, 2025-1.}

            \ecvitem*{}{2023. \textbf{Mercadeo predictivo}. Maestría en Mercadeo. Periodos (1 semestre): 2023-2. Diseño original del curso.}

        \ecvitem{Universidad de los Andes -- Educación ejecutiva}{}
            \ecvitem*{}{2022 - 2023. \textbf{Escuela de Finanzas Corficolombiana (cohortes 2 y 3)}.}

            \ecvitem*{}{2016 - 2017. \textbf{Análisis de datos, Banco Popular}.}

        \ecvitem{Docencia internacional}{}
            \ecvitem*{}{2026. \textbf{Blockchain y criptoactivos (con Aarón Olmos)}. IESA, Caracas, Venezuela. Programa conjunto IESA--Uniandes. Fechas: enero--marzo 2026.}

            \ecvitem*{}{2024. \textbf{Taller III: Finanzas aplicadas}. Global MBA, Universidad Americana, Asunción, Paraguay. Fechas: marzo--junio 2024.}

        \ecvitem{Universidad del Rosario}{}
            \ecvitem*{}{2010 - 2013. \textbf{Econometría básica}. Pregrado en Economía. Periodos (6 semestres): 2010-2, 2011-1, 2011-2, 2012-1, 2012-2, 2013-1.}

            \ecvitem*{}{2010 - 2013. \textbf{Econometría intermedia}. Pregrado en Economía. Periodos (6 semestres): 2010-2, 2011-1, 2011-2, 2012-1, 2012-2, 2013-1.}

            \ecvitem*{}{2013. \textbf{Introducción a la economía}. Pregrado en Economía. Periodos (1 semestre): 2013-1.}

            \ecvitem*{}{2013. \textbf{Curso básico de Stata}. Pregrado en Economía. Periodos (1 semestre): 2013-1.}

            \ecvitem*{}{2004. \textbf{Macroeconomía II (economía abierta)}. Pregrado en Economía. Periodos (2 semestres): 2004-1, 2004-2.}

            \ecvitem*{}{2003. \textbf{Economía para no economistas}. Pregrado en Economía. Periodos (1 semestre): 2003-2.}

            \ecvitem*{}{2002 - 2003. \textbf{Macroeconomía I (economía cerrada)}. Pregrado en Economía.}

            \ecvitem*{}{2002. \textbf{Introducción a la microeconomía}. Pregrado en Economía y Ciencia Política. Periodos (1 semestre): 2002-1.}

            \ecvitem*{}{2002. \textbf{Economía de la regulación (asistente de profesor)}. Pregrado en Economía. Periodos (1 semestre): 2002-2.}

        \ecvitem{Australian National University}{}
            \ecvitem*{}{2007. \textbf{Tutor de economía, JABAL Centre}. Australian National University.}

            \ecvitem*{}{\footnotesize{\emph{int = periodo intersemestral.}}}

%---------------------------------------------------------------------
    \cline{2-2}
    \ecvsection{Dirección de tesis}

        \ecvitem{Doctorado en Administración}{}
            \ecvitem*{}{2026 - . \textbf{Nicolás Felipe Suárez}. Director (codirector: Daniel Mantilla). En curso.}

            \ecvitem*{}{2025 - . \textbf{John Jairo Gálvis}. \emph{La brecha universidad--mercado en programas de administración de empresas: una explicación multinivel desde currículo, mercado laboral e instituciones}. Propuesta sustentada en 2026. En curso.}

        \ecvitem{Maestría}{}
            \ecvitem*{}{2023. \textbf{José David Pinzón Melo}. \emph{From bonds to prosperity: echoes of yield spread on GDP growth}. International Master in Finance, NOVA--Uniandes. Supervisor Uniandes.}

            \ecvitem*{}{2022. \textbf{Elkin Julián Salcedo Sánchez}. \emph{Structural and speculative factors behind the S\&P 500 futures real estate market}. International Master in Finance, NOVA--Uniandes. Cosupervisión con Melissa Prado.}

            \ecvitem*{}{2020. \textbf{Juan Diego Londoño Herrera}. \emph{Efecto de la regulación bancaria sobre el crecimiento económico: costos y política macroprudencial de largo plazo}. Maestría en Economía (PEG), Facultad de Economía.}

        \ecvitem{Pregrado}{}
            \ecvitem*{}{2019. \textbf{Isabella Arias Ortiz y Laura Cristina Herrán Sánchez}. \emph{Tipo de financiamiento y desempeño según ventas en microestablecimientos colombianos}. Memoria de grado, Facultad de Economía.}

            \ecvitem*{}{2017. \textbf{Daniel Mora Ardila}. \emph{Sobrevaloración en el mercado de vivienda nueva en las principales ciudades de Colombia}. Memoria de grado, Facultad de Economía.}

        \ecvitem{Jurado}{}
            \ecvitem*{}{Jurado o lector de seis trabajos de grado de pregrado (2013--2017) y de 23 de maestría (2013--2023), entre ellos del International Master in Finance NOVA--Uniandes y del MBA.}

%---------------------------------------------------------------------
    \cline{2-2}
    \ecvsection{Premios y distinciones}

        \ecvitem*{}{\textbf{Banco de la República. (Banco Central de Colombia)}. Beca para estudios de doctorado en economía en el exterior. Junio 2004.}

        \ecvitem*{}{\textbf{COLCIENCIAS}. Joven investigador. Otorgada para llevar a cabo actividades de investigador en la Universidad del Rosario. Bogota, Colombia. Octubre 2000.}

%---------------------------------------------------------------------
    \cline{2-2}
    \ecvsection{Evaluador}

        \ecvitem{Revistas}{}
            \ecvitem*{}{Academia Revista Latinoamericana de Administración, Economic Change and Restructuring, Economic Papers, Economic Systems, Emerging Markets Finance and Trade, Empirical Economics, Energy Economics, Energy Journal, Ensayos sobre Política Económica, Estudios Gerenciales, International Review of Economics and Finance, Journal of Business Research, Lecturas de Economía, Managerial Finance, Revista de Economía del Rosario, Sociedad y Economía, Stata Journal.}

        \ecvitem{Casos}{}
            \ecvitem*{}{North American Case Research Association (NACRA).}

        \ecvitem{Proyectos}{}
            \ecvitem*{}{Swiss National Science Foundation (SNSF), COLCIENCIAS, Banco de la República, Fundación para la Promoción de la Investigación y la Tecnología (FPIT, Banco de la República), y convocatorias y publicaciones de universidades colombianas: Antioquia, Atlántico, Cartagena, Cauca, EAN, EIA, ESAP, Guajira, Javeriana, La Salle, Medellín, Militar, Norte, Pedagógica y Tecnológica (UPTC), Rosario, Francisco de Paula Santander (UFPS), Industrial de Santander (UIS), Sucre, Surcolombiana.}

        \ecvitem{Premios}{}
            \ecvitem*{}{Jurado, Premio Otto de Greiff (2015, 2025).}

%---------------------------------------------------------------------
    \cline{2-2}
    \ecvsection{Idiomas}
        \ecvmothertongue{Español}
        \ecvlanguageheader{*}
        \ecvlastlanguage{Inglés}{\ecvCOne}{\ecvCOne}{\ecvCOne}{\ecvCOne}{\ecvCOne}
        \ecvlanguagefooter{(*)}

%---------------------------------------------------------------------
    \cline{2-2}
    \ecvsection{Otras habilidades}
        \ecvitem{\footnotesize{\emph{Durante el período de estudio y trabajo como investigador académico he adquirido habilidad en el uso de los siguientes sistemas}}}

        \ecvitem*{}{Microsoft Office (Word, Excel, PowerPoint, Access, Outlook).}

        \ecvitem*{}{\LaTeX}

        \ecvitem*{}{Programas de econometría: Stata, Eviews, R.}

        \ecvitem*{}{Visualización de datos: Power BI.}

        \ecvitem*{}{Programas de estimacion de eficiencia técnica: EMS, Warwick DEA Software.}

        \ecvitem*{}{Programas de diseño de material para Internet: Dreamweaver, Fireworks.}

%---------------------------------------------------------------------
    \cline{2-2}
    \ecvsection{Otras actividades}
        \ecvitem{\footnotesize{\emph{Otras actividades desarrolladas conjuntamente con mis actividades de economista}}}

        \ecvitem*{}{Editor Revista de Economía del Rosario - Departamento de Economía, Universidad del Rosario. Enero 2010 - Diciembre 2023}

        \ecvitem*{}{Asistente Revista de Economía del Rosario - Departamento de Economía, Universidad del Rosario. Enero 2003 - Diciembre 2004}

        \ecvitem*{}{Editor Borradores de Investigación - Departamento de Economía, Universidad del Rosario. Enero 2003 - Diciembre 2004}

        \ecvitem*{}{Construcción, edición y manejo de sitio internet del Departamento de Economía de la Universidad del Rosario. Mayo 2000 - Diciembre 2004 \emph{http://economia.urosario.edu.co}.}

        \ecvitem*{}{Construcción, edición y manejo de sitio internet del Departamento de Economía de la Universidad de Antioquia. Enero 1999 - Abril 2000 \emph{http://economia.udea.edu.co/economia}.}

        \ecvitem*{}{Editor Revista Oikos. (Revista de economía editada por los estudiantes de pregrado - Universidad de Antioquia) Enero 1999 - June 2000}

%---------------------------------------------------------------------
    \cline{2-2}
    \ecvsection{Referencias}
        \ecvitem{\footnotesize{\emph{Personal and professional references can be asked to}}}

%        \ecvitem*{}{Carlos Pombo. Facultad de Administración. Universidad de los Andes, Bogotá. \emph{cpombo@uniandes.edu.co}. Telephone: (+ 57 1) 339 4949.}

%        \ecvitem*{}{Luis Gutiérrez. Facultad de Economía. Universidad del Rosario, Bogotá. \emph{luis.gutierrez@urosario.edu.co}. Telephone: (+ 57 1) 2970200.}

%        \ecvitem*{}{Juan Oviedo. Facultad de Economía. Universidad del Rosario, Bogotá. \emph{juan.guataqui@urosario.edu.co}. Telephone: (+ 57 1) 2970200.}

%        \ecvitem*{}{Peter Warr. RSPAS-Economics. Australian National University, Canberra. \emph{peter.warr@anu.edu.au}. Telephone: (+ 61 2) 612 52682.}
%
%        \ecvitem*{}{Martin Richardson. Economics faculty. Australian National University, Canberra. \emph{martin.richardson@anu.edu.au}. Telephone: (+ 61 2) 612 53582.}
%
%        \ecvitem*{}{John Soderbaum. \emph{j.soderbaum@aciltasman.com.au}. Telephone: (+ 61 2) 610 38200..}
%
%        \ecvitem*{}{Lynne Gallagher. \emph{l.gallagher@aciltasman.com.au}. Telephone: 0412255704.}
%
%        \ecvitem*{}{Alan Smart. ACIL-Tasman (Canberra) \emph{a.smart@aciltasman.com.au}. Telephone: (+ 61 2) 610 38200.}

\end{europecv}

\end{document}
